% =============================================================
%  ABSTRACT
% =============================================================
\chapter*{Abstract}
\addcontentsline{toc}{chapter}{Abstract}

\begingroup
\setlength{\parskip}{0pt}
\setstretch{1.2}

\noindent Clustering on social datasets assembled from heterogeneous sources is
prone to recovering data-collection-source structure before phenomenon-level
structure, while standard internal validity indices can reward such degenerate
partitions as ``tight'' clusterings. Reliable analysis on this kind of input
therefore depends on the joint choice of imputation method, clustering
algorithm, and evaluation metric. This study constructs and systematically
evaluates a factorial benchmark of these three choices on 1,617 incomplete,
newspaper-derived Bangladeshi suicide cases (KNN and MissForest variants impute
all 1,617 rows; the MICE variants retain 1,613 after dropping four rows whose
near-empty predictor profile prevented convergent imputation) --- the largest
publicly assembled record
of its kind for a country that lacks any centralised national suicide
registry. Three imputation methods (MICE, KNN, MissForest) and four clustering
algorithms (K-Means, Agglomerative, DBSCAN, Spectral) are evaluated under two
predictor strategies (context-aware, no-context), producing 32 pipelines plus
two encoded baselines, each scored by a six-component composite (Silhouette,
Calinski--Harabasz, Davies--Bouldin, Dunn, Xie--Beni, and a study-specific
balance penalty designed to discourage near-degenerate cluster assignments).
The best configuration --- Agglomerative + MissForest, no-context --- achieves
a composite of 27.30/32 and a two-cluster solution that aligns almost
perfectly with the dataset's 2020 Kaggle vs.\ 2017--2019 + 2021--2025 manual
cohort boundary. The cluster boundary is therefore partly confounded
with collection source: any signal it captures is intertwined with reporting
era, completeness, and source heterogeneity, and the results should be read as
exploratory subgroup structure rather than discovery of distinct
epidemiological categories. Subject to this caveat, the framework
surfaces stable, interpretable structure from severely incomplete newspaper
data and exposes source-cohort artefacts that simpler single-method workflows
would conceal.

\vspace{0.15in}
\noindent Keywords: suicide patterns; missing data; multiple imputation;
MissForest; clustering; cluster validity; Bangladesh.

\endgroup
\clearpage
